% texbook.py 注入的页眉/页脚样式（生成物，勿手改）
\usepackage{fancyhdr}
\pagestyle{fancy}
\fancyhf{}
\fancyhead[L]{\small\color{gray}Python 科学计算（新版教程）}
\fancyhead[R]{\small\color{gray}\thepage}
\renewcommand{\headrulewidth}{0.4pt}

% ===== 全局排版：行距/表格/图片/断行 =====
\usepackage{setspace}
\setstretch{1.15}
\renewcommand{\arraystretch}{1.18}
\setlength{\tabcolsep}{4pt}
\emergencystretch=3em
\clubpenalty=4000
\widowpenalty=4000
\RequirePackage{graphicx}

% ===== 图片：限制高度并只在页顶/独立页排版，避免被页面下沿裁切 =====
\usepackage{float}
\floatplacement{figure}{tp}
\setkeys{Gin}{width=\textwidth,height=0.70\textheight,keepaspectratio}
\setlength{\intextsep}{10pt}
\setlength{\textfloatsep}{12pt}

% ===== 代码块：长行自动折行（长注释、长命令不再溢出到页边）=====
\usepackage{fvextra}
\fvset{breaklines=true,breakanywhere=true,breakindent=0pt,breaksymbolleft={},breaksymbolright={},fontsize=\small}
\DefineVerbatimEnvironment{Highlighting}{Verbatim}{commandchars=\\\{\},breaklines=true,breakanywhere=true,breakindent=0pt,fontsize=\small}

% ===== 正文断行：长 URL 与长等宽串允许断行 =====
\usepackage{xurl}
\usepackage[htt]{hyphenat}
\sloppy
\hbadness=10000

% 长表格用 \small 字体，减少列溢出
\usepackage{longtable}
\usepackage{etoolbox}
\AtBeginEnvironment{longtable}{\small}
